\section{Soft triangles are the source of hardness}\label{sec:soft}
The remaining sections locate where the difficulty of CEaMP sits. For zero-excess solutions the interaction between packed $P_3$'s is as simple as a constraint system can be unless the packing contains one specific configuration, which we call a \emph{soft triangle}. Without it, zero excess is polynomial for every $c$ and $w$ (Corollary~\ref{cor:sfpoly}), every lock core has at most two anchors (Corollary~\ref{cor:sfcore}), and a far-lift model of excess is fixed-parameter tractable (Theorem~\ref{thm:U}); on forests $I(H)$ the full problem is fixed-parameter tractable in $\ell$ conditionally on one bounded-certificate hypothesis (Theorem~\ref{thm:F}). With soft triangles, the natural branching scheme provably has unbounded branching (Theorem~\ref{thm:CE}). The main question of Section~\ref{sec:open} stays open; the new, more precise status is stated in Section~\ref{sec:softstatus}.

\paragraph{Conventions and labels.}
States are as in Appendix~\ref{app:tunnels}: a member $(a,b,c)$ with centre $b$ is in state $1,2,3$ if its single edited pair is $ab,bc,ac$. Put $t_{ab}=(0,1,1)$, $t_{bc}=(1,0,1)$, $t_{ac}=(0,0,1)$ for the states $(1,2,3)$: $t_e(s)=1$ iff the pair $e$ of the member is inside one cluster after the edit. $\Sigma=\Sigma(G,H)\subseteq\{1,2,3\}^H$ is the set of zero-excess (ZE) state vectors; by Lemma~\ref{lem:decomp}, $s\in\Sigma$ iff no vertex triple is an induced $P_3$ of $G$ after the edits $s$. A \emph{majority operation} is $m:\{1,2,3\}^3\to\{1,2,3\}$ with $m(a,a,b)=m(a,b,a)=m(b,a,a)=a$; applied to vectors it acts coordinatewise. Every statement is marked: \emph{proved} (a proof is given, in Appendix~\ref{app:soft} when long), \emph{by computer} (a finite case check that is part of the proof, or of an example, and that was carried out by two independent implementations), or \emph{evidence} (computer experiments that are not part of any proof). The statements of this section were also re-derived independently (new code, exhaustive enumeration of all set partitions for $n\le 11$); this is evidence and not a substitute for the proofs.

\subsection{Soft triangles and the structure of zero-excess solutions}
\begin{definition}[Soft triangle]\label{def:softtri}
A vertex triple $\{x,y,z\}$ is a \emph{soft triangle} of $(G,H)$ if its three pairs are soft and are owned by three pairwise distinct members. $(G,H)$ is \emph{soft-triangle-free} if there is none. (This depends on $H$ only, not on the edges of $G$.)
\end{definition}
\begin{lemma}[Soft triangles and $I(H)$]\label{lem:tri}
(a) The owners of a soft triangle form a triangle of $I(H)$. (b) If $c\le2$, every triangle of $I(H)$ consists of the owners of a soft triangle. (c) If $I(H)$ is triangle-free, in particular if $I(H)$ is a forest ($w\le1$), there is no soft triangle. For $c\ge3$ a triangle of $I(H)$ need not come from a soft triangle (three members through one vertex).
\end{lemma}
\begin{proof}
(a) Two pairs of a triple share a vertex, so their owners intersect. (b) Let $A,B,C$ be pairwise adjacent. If $A\cap B$, $B\cap C$ and $A\cap C$ shared a vertex, it would lie in three members, contradicting $c\le2$; so the three intersection vertices $x\in A\cap B$, $y\in B\cap C$, $z\in A\cap C$ are distinct, $\{x,y\}\subseteq V(B)$, $\{y,z\}\subseteq V(C)$, $\{x,z\}\subseteq V(A)$ are soft pairs with three distinct owners. (c) follows from (a).
\end{proof}

\begin{lemma}[Constraint structure]\label{lem:sfstruct}
Let $H$ be any packing and classify a vertex triple $\tau$ by the owners of its soft pairs.
\begin{enumerate}
\item[(i)] No soft pair: a \emph{static} condition on $G$ (if $\tau$ is an induced $P_3$ of $G$, then $\Sigma=\emptyset$).
\item[(ii)] One soft pair, owner $A$: a \emph{unary} constraint on $s_A$.
\item[(iii)] Two soft pairs, from distinct members $A,B$: these share a vertex $q\in V(A)\cap V(B)$, $\tau=\{p,q,r\}$ with $p\in V(A)\setminus\{q\}$, $r\in V(B)\setminus\{q\}$, and $pr$ is hard; a \emph{binary} constraint on $(s_A,s_B)$, namely $t_{pq}(s_A)+t_{qr}(s_B)+[pr\in E]\ne2$.
\item[(iv)] All soft pairs from one member: $\tau=V(A)$, no restriction on $\Sigma$ beyond the states.
\item[(v)] Three soft pairs from three distinct members: a soft triangle, a \emph{ternary} constraint.
\end{enumerate}
Hence if $(G,H)$ is soft-triangle-free, $\Sigma$ is empty or the set of vectors satisfying unary lists $\Lambda_A\subseteq\{1,2,3\}$ and, for every pair $A,B$ of adjacent members sharing $q$, the binary relation $R_{AB}\subseteq\{1,2,3\}^2$ that is the conjunction of the four triples $\{p,q,r\}$.
\end{lemma}
\begin{proof}
Two soft pairs of one triple share a vertex. If their owners coincide, the triple is $V(A)$, which is case (iv). If they differ, the shared vertex lies in both members, so it is the unique common vertex $q$; the third pair $pr$ is not soft, since an owner $A$ or $B$ would then contain two vertices of the other member (forcing $|V(A)\cap V(B)|\ge2$) and a third owner gives (v). The constraint of (iii) is ``not exactly two edges''. The final sentence is the classification, as $R_{AB}$ depends only on $A,B$ and the bits $[pr\in E]$ (the cross pairs).
\end{proof}

\begin{theorem}[Majority closure]\label{thm:sfA}
If $(G,H)$ is soft-triangle-free, then $\Sigma(G,H)$ is closed under \emph{every} majority operation $m$ applied coordinatewise, with the same $m$ at all coordinates: for $u,v,w\in\Sigma$, $m(u,v,w)\in\Sigma$. No pair-maximality is needed, and $c$ and $w$ are arbitrary.
\end{theorem}
\emph{Proved; the classification of the binary relations $R_{AB}$ is by computer} (Appendix~\ref{app:softA}). The relation $R_{AB}$ depends on the roles of $q$ in $A$ and in $B$ and on the four cross bits: $144$ configurations, giving exactly $40$ distinct relations, all of the shapes empty, singleton, $2$-element set, a full row or column through $\{1,2\}$, a $3$-element subset of $\{1,2\}^2$, or a ``cross'' (row $\cup$ column through $\{1,2\}$); each shape is closed under every majority by a short argument. The classification was reproduced by an independent construction from $576$ real graphs, and all $729$ majority operations were checked against all relations. The point of the theorem is that single triples are not closed under any uniform majority; only the conjunction of the four triples at a shared vertex is.

\begin{corollary}[Zero excess without soft triangles is polynomial]\label{cor:sfpoly}
For soft-triangle-free $(G,H)$, whether $(G,H,0)$ is a YES instance is decidable in polynomial time. No bound on $c$ or $w$ and no tree decomposition is needed.
\end{corollary}
\begin{proof}
By Lemma~\ref{lem:sfstruct}, $(G,H,0)$ is YES iff there is no static violation and the binary CSP with $|H|$ variables of domain $\{1,2,3\}$, unary lists $\Lambda_A$ and relations $R_{AB}$ is satisfiable. By Theorem~\ref{thm:sfA} every constraint is closed under a majority operation. By the theorem of Jeavons, Cohen and Cooper~\cite{JCC98}, a set of constraints closed under a near-unanimity operation (here a majority, of arity $3$) is decided by enforcing strong consistency of a fixed level, which for binary constraints and a majority is path consistency; this takes time polynomial in $|H|$. Static triples are tested in $O(n^3)$.
\end{proof}
Evidence: an independent path-consistency solver built from the vertex triples agreed with exhaustive enumeration on $20{,}920$ random soft-triangle-free instances ($n\le10$, $c\le4$, feasible and infeasible) and $3{,}400$ tree-like ones, and every one of them had $\Sigma$ closed under all $729$ majority operations; on instances \emph{with} soft triangles the same code finds $451$ failures of closure.

\begin{remark}[No conflict with the hardness of $\ell=0$]\label{rem:sfLPS}
Corollary~\ref{cor:sfpoly} does not contradict the NP-hardness of $\ell=0$ for $c\le23$~\cite{LPS23}: the variable gadget of~\cite[Sect.~4.1.1]{LPS23} (three cliques of five vertices over $\mathbb F_5$, one member for each arithmetic progression, $25$ members, $c=5$) packs every pair between consecutive cliques and a computer count finds $100$ soft triangles in it (evidence; any triple with one vertex in each clique that is not an arithmetic progression is one). The hardness proof therefore lives on ternary constraints. Likewise Example~\ref{ex:maj} ($c=2$) has four soft triangles and is not closed under the operation $m_3$, and the flip gadget of Example~\ref{ex:gstar} has soft triangles. At $c\le2$ a soft triangle is a triangle of $I(H)$ (Lemma~\ref{lem:tri}), hence $w\ge2$, so for $w\le1$ the hardness of $\ell=0$ cannot occur (this is also clear from Theorem~\ref{thm:l0}).
\end{remark}

\begin{corollary}[Lock cores have at most two anchors]\label{cor:sfcore}
If $(G,H)$ is soft-triangle-free, every core (minimal non-face of anchors, Appendix~\ref{app:tunnels}) has at most two anchors. In particular this holds for every $H$ with $c\le2$ and $I(H)$ triangle-free, whatever $\tw(I(H))$.
\end{corollary}
\begin{proof}
Theorem~\ref{thm:sfA} gives closure under a majority operation, and Lemma~\ref{lem:tmaj} applies.
\end{proof}
This settles, for these instances, the question whether the hypothesis of Lemma~\ref{lem:tmaj} holds: it holds exactly when there is no soft triangle in the sense that matters. The statement that the specific operation $m_3$ closes $\Sigma$ for $c\le2$ is false with triangles (Example~\ref{ex:maj}) and true without soft triangles.

\begin{corollary}[Tunnel releases are fixed-parameter tractable without soft triangles]\label{cor:sfTR}
In the tunnel-release model of Appendix~\ref{app:tunnels}, for a soft-triangle-free packing the minimum number of cuts is computable in time $2^{\ell}n^{O(1)}$, with no dependence on $c$ or $w$.
\end{corollary}
\begin{proof}
A set $S$ of tunnels is lockable iff the anchors at the ends of its tunnels in $V(H)$ form a face of $G$ with all tunnels removed (Lemma~\ref{lem:tinactive}). By Corollary~\ref{cor:sfcore} a non-face contains a non-face of at most two anchors, hence of at most two tunnels, so every core has at most two tunnels. Proposition~\ref{prop:tTR}(iv) with $d=2$ applies, with the $\ell=0$ test of Corollary~\ref{cor:sfpoly}.
\end{proof}

\begin{remark}[With triangles at $c\le2$: evidence only]\label{rem:tri2}
For $c\le2$ with triangles in $I(H)$ the set $\Sigma$ can fail to be closed under a given majority (Example~\ref{ex:maj}). Whether it is always \emph{$2$-decomposable} (equal to the set of vectors whose binary projections lie in the binary projections of $\Sigma$), which is what Lemma~\ref{lem:tmaj} actually needs, is open. Evidence: all tight triangle blocks (three members on six vertices, all role orders and hard bits; $3{,}712$ tight configurations, $163$ distinct $\Sigma$), all tight $K_4$ gadgets ($521$ distinct $\Sigma$), $3{,}000$ random tight instances on the triangular prism and random instances with $c\le2$ all gave $2$-decomposable $\Sigma$ and no core with three anchors, although $24$ of the $163$ sets $\Sigma$ of the blocks and $56$ of the $521$ of the $K_4$ gadgets are not closed under all majorities. In an independent re-derivation, among random instances with soft triangles the non-$2$-decomposable ones ($34$ in the runs that recorded $c$) all had $c=3$, and instances with $c=2$ that are not closed under all majorities did occur ($33$).
\end{remark}

\subsection{Far hard pairs: what they can and cannot do}\label{sec:softfar}
Call a hard pair $\{p,r\}$ a \emph{cross pair} if $p\in V(A)\setminus\{q\}$ and $r\in V(B)\setminus\{q\}$ for two members $A\ne B$ with $V(A)\cap V(B)=\{q\}$ (these are the cross bits of Lemma~\ref{lem:sfstruct}); every other hard pair is \emph{far}. A far pair never lies in a binary triple.

\begin{lemma}[U-only identity and locality of hard edits]\label{lem:sfU}
Call a solution \emph{U-only} if every member receives exactly one edit, so that its other edits form a set $U$ of hard pairs. For every packing, the least size of a U-only solution is $|H|+\min\{|U|:\Sigma(G\sd U,H)\ne\emptyset\}$. For soft-triangle-free $(G,H)$, $\Sigma(G\sd U,H)$ differs from $\Sigma(G,H)$ only through (a) the static triples (these gate non-emptiness: an all-hard induced $P_3$ created or destroyed by $U$ empties or un-empties $\Sigma$ globally), (b) the lists $\Lambda_A$ of the members $A$ that contain an end of a pair of $U$, and (c) the relations $R_{AB}$ of adjacent members $A,B$ for which $U$ contains a cross pair of $(A,B)$. In particular a far pair changes only static triples and the lists at the members through its two ends.
\end{lemma}
\begin{proof}
$U$ avoids soft pairs, so every member is still an induced $P_3$ of $G\sd U$, with the same pairs, and a vector is the single-edit part of a solution of $G$ with hard edits $U$ iff it lies in $\Sigma(G\sd U,H)$. A hard pair $e$ enters only the triples containing it; by Lemma~\ref{lem:sfstruct} these are static, unary on a member through an end of $e$, or (if $e$ is a cross pair of $(A,B)$) in a binary triple of $R_{AB}$ or in a list of $A$ or $B$.
\end{proof}
The model is a sub-model: one edit per member is not always optimal, as in Remark~\ref{rem:domain}; for instance $n=5$, $E=\{02,03,12,23,24\}$, $H=\{(1,2,4)\}$ has $\OPT=|H|+1$ but least U-only value $|H|+2$ (by enumeration of all partitions).

\begin{lemma}[Strictness]\label{lem:sfS}
\emph{(by computer; finite check)} Let $A,B$ be adjacent members of a packing, $R$ the relation $R_{AB}$, and $\Lambda_A,\Lambda_B$ the lists that the cross bits of $(A,B)$ impose on $A$ and on $B$. For $X\subseteq\{1,2,3\}$ put $\mathrm{img}(X)=\{s_A:\exists s_B\in X\cap\Lambda_B,\ (s_A,s_B)\in R\}\cap\Lambda_A$. If $X\subseteq Y\subseteq\{1,2,3\}$ and $\mathrm{img}(X)\ne\mathrm{img}(Y)$, then $|\mathrm{img}(X)|\le1$.
\end{lemma}
Across an edge of $I(H)$, a restriction of the neighbour's list that really changes the image leaves at most one state; locks propagate as unit propagation along edges of $I(H)$. This is a one-edge statement, and was checked for $144$ configurations and $27$ pairs $X\subseteq Y$ ($803$ strict cases, no violation); an independent check on $576$ real graphs agrees. It is used only through Lemma~\ref{lem:sfH2}.

\begin{lemma}[Helly number two for unit lists]\label{lem:sfH2}
Let $\Sigma\subseteq D^V$ be closed under some majority operation (for any finite domain $D$) and let $J$ be a family of unit constraints $x_{v_j}=s_j$. If $\Sigma$ and $J$ are jointly infeasible, then $\Sigma$ together with some at most two members of $J$ is infeasible.
\end{lemma}
\begin{proof}
Take a minimal infeasible $J$ with $|J|=k\ge3$ and, for each $j$, $t_j\in\Sigma$ satisfying $J\setminus\{j\}$. For three indices $a,b,c$ put $u=m(t_a,t_b,t_c)\in\Sigma$. For $j\notin\{a,b,c\}$ all three tuples have value $s_j$ at $v_j$; at $v_a$ the tuples $t_b,t_c$ both have value $s_a$, so $u$ has $s_a$ there (and symmetrically for $b,c$). Hence $u$ satisfies all of $J$, a contradiction. (If two constraints concern one coordinate with different values, the same argument applies.)
\end{proof}

\begin{observation}[Lists of an outside vertex]\label{obs:lists}
Let $A=(a,b,c)$ and let $t\notin V(A)$ be such that the three pairs $\{x,t\}$, $x\in V(A)$, are hard; put $S=N(t)\cap V(A)$. The three triples $\{a,b,t\},\{b,c,t\},\{a,c,t\}$ allow exactly the following states of $A$:
\begin{center}
\begin{tabular}{c|cccccccc}
$S$ & $\emptyset$ & $\{a\}$ & $\{b\}$ & $\{c\}$ & $\{a,b\}$ & $\{a,c\}$ & $\{b,c\}$ & $\{a,b,c\}$\\\hline
allowed in $\{1,2,3\}$ & $\{1,2,3\}$ & $\{1\}$ & $\emptyset$ & $\{2\}$ & $\{2\}$ & $\emptyset$ & $\{1\}$ & $\{3\}$\\
allowed in $\{1,\dots,5\}$ & $\{1,\dots,5\}$ & $\{1,4\}$ & $\{4,5\}$ & $\{2,4\}$ & $\{2\}$ & $\{5\}$ & $\{1\}$ & $\{3\}$
\end{tabular}
\end{center}
(Direct check; the second row is for the five states of Lemma~\ref{lem:five}: $S$ must be empty or a single part of the cluster partition of $\{a,b,c\}$ after the edit.) So a zero-excess list of an outside vertex is trivial, empty, or a \emph{unit} list; there are no two-element lists. This is why forked locks never occur for zero excess.
\end{observation}

\begin{theorem}[Far-lift zero-excess model]\label{thm:U}
Let $(G,H)$ be soft-triangle-free, with $c$ and $w$ arbitrary. It is decidable in time $6^{\ell}n^{O(1)}$ whether there is a set $W$ of at most $\ell$ far pairs with $\Sigma(G\sd W,H)\ne\emptyset$ (the \emph{far-lift model}). A YES answer gives a solution of CEaMP of size $|H|+|W|$.
\end{theorem}
\emph{Proved} (Appendix~\ref{app:softU}). Algorithm: (1) if $G'=G\sd W$ has an all-hard induced $P_3$, branch on its far pairs (none: NO); (2) otherwise accept if $\Sigma(G')\ne\emptyset$ (Corollary~\ref{cor:sfpoly}); (3) delete all \emph{liftable} unit lists (those of Observation~\ref{obs:lists} for which at least one of the three pairs is far); if the rest is already infeasible, answer NO; (4) otherwise find a minimal infeasible family of liftable lists, which has at most two members by Lemma~\ref{lem:sfH2} (the rest of $\Sigma$ is closed under all majorities by Theorem~\ref{thm:sfA}), and branch on the at most $3$ far pairs of each: at most $6$ branches, depth at most $\ell$. A far pair changes only liftable lists (Lemma~\ref{lem:sfU}), so some far pair of the family must be edited. \emph{Evidence}: the algorithm agreed with exhaustive search on $1{,}420$ instances (and, in an independent re-implementation, on $1{,}050$ more; the largest minimal family had two lists).

The model is restricted (no double edits, no edits of cross pairs): it is an upper bound on the excess, exact when some optimum uses neither. It contains the tunnel-release model of Corollary~\ref{cor:sfTR} and removes the need for tunnels: \emph{arbitrary} far hard edits, including those that close static violations, are handled. It also shows what the cores of far edits look like: a pair of unit lists, a vertex-cover-like structure whose branching factor depends neither on $n$, nor $c$, nor $w$.

\subsection{Five states, an exact model, and the abstract models}\label{sec:softmodel}
By Lemma~\ref{lem:five} the variables of the dynamic programs of Section~\ref{sec:csp} range over five states: besides $1,2,3$, the state $4$ (edits $ab,bc$; extra cost $1$) and $5$ (edits $ab,bc,ac$; extra cost $2$). For a graph $G'$ with the same packing let $\mathrm{Ex}(G')$ be the least value of $|F|+\sum_i\mathrm{extra}(x_i)$ over sets $F$ of cross pairs and state vectors $x\in\{1,\dots,5\}^H$ such that $G'\sd F\sd(\text{soft edits of }x)$ is a cluster graph ($\infty$ if none); extra $=0,0,0,1,2$ for the states $1,\dots,5$.
\begin{lemma}[Exactness]\label{lem:sfE}
For every packing, $\OPT(G)=|H|+\min_{W}\bigl(|W|+\mathrm{Ex}(G\sd W)\bigr)$, where $W$ ranges over sets of far pairs.
\end{lemma}
\begin{proof}
Split a solution as in Lemma~\ref{lem:decomp} into its soft edits (one non-empty edit set per member, which by Lemma~\ref{lem:five} is a state; its cost is $1+\mathrm{extra}$), its edits of cross pairs $F$, and its edits of far pairs $W$. Conversely every $(W,F,x)$ is a solution of that size.
\end{proof}
Hence the \emph{far} edits are the only edits whose position is not controlled by the dynamic program, which is what Section~\ref{sec:nl} found.

\paragraph{Breaks.} Let $\mathrm{RM2}$ be the following abstract model on a forest: a set of \emph{conflicts} (pairs of nodes, with the path joining them) and two kinds of operations, a \emph{break} at a node (cost $1$; it hits every conflict whose path contains the node) and a \emph{tunnel}, a link in a graph of maximum degree $d$ (cost $1$; it hits every conflict that has an end at one of its two nodes). We want at most $L$ operations hitting all conflicts.
\begin{proposition}[RM2 is fixed-parameter tractable]\label{prop:RM2}
RM2 is solvable in time $(2d+1)^{L}n^{O(1)}$ by a search tree.
\end{proposition}
\begin{proof}
Root every tree and let $Q$ be an unhit conflict whose top node $v$ is deepest. If an optimal solution hits $Q$ by a break at a node $w\in Q$, replace it by a break at $v$: every other unhit conflict $R$ through $w$ has its top at depth at most $\mathrm{depth}(v)$ and contains $w$, hence contains the whole path from $w$ to its top, which passes through $v$. If $Q$ is hit by a tunnel, the tunnel has an end at an end of $Q$ ($2d$ choices). So the search tree branches $2d+1$ ways to depth $L$.
\end{proof}
(Evidence: agreement with exhaustive search on $5{,}000$ + $2{,}100$ random instances in two independent implementations.) RM2 was a plausible abstraction of CEaMP on forests; it is not a faithful one.
\begin{proposition}[RM2 is not faithful]\label{prop:hub}
Let $n=11$ (vertices $0,\dots,10$), let $G$ have the twelve edges
\[01,\ 02,\ 04,\ 05,\ 06,\ 13,\ 28,\ \{2,10\},\ 34,\ 37,\ 38,\ 78\]
and let $H=\{M_0,M_1,M_2,M_3\}$ with $M_0=(2,0,1)$, $M_1=(4,3,1)$, $M_2=(5,0,6)$, $M_3=(2,8,7)$. This instance is valid, pair-maximal, $I(H)$ is a star with centre $M_0$, $c=2$ and there is no soft triangle. Then $\OPT=|H|+3$ (enumeration of all set partitions). The members $M_1$ and $M_2$ cannot be in any of the states $1,2,3$ in a solution without hard edits (an outside vertex gives them the list $\{4,5\}$ of the second row of Observation~\ref{obs:lists}), while $M_0,M_3$ can; each of $M_1,M_2$ is repaired alone by one break (state $4$), so the break model predicts excess $2$, but the true excess is $3$: two breaks next to a common member put the cluster-views $(0,0)$ at two vertices of $M_0$, which then cannot be in a state of $\{1,2,3\}$ either. The cheapest vectors without hard edits are $(1,5,4,1)$ and $(4,4,4,1)$, both of cost $7$.
\end{proposition}
The proposition (checked by exhaustive enumeration, and by an independent subset dynamic program) says that a break is not a node-hit: it locks its neighbours. Consequently the abstract model of ``hitting conflict paths by breaks and tunnels'' cannot be used to decide CEaMP on forests, and the difficulty of the interior of a chain, where breaks are placed, stays inside the dynamic program (as the Core-Bound analysis below also shows).

\paragraph{The full problem on forests, conditionally.}
Let $I(H)$ be a forest ($w\le1$, so $c\le2$, $D=2c\le4$, no soft triangle). The \emph{sites} of $(G,H)$ are the members and the edges of $I(H)$; a triple that contains a far pair has as sites the owners of its soft pairs and the edges at which its cross pairs lie. For a set $C$ of sites let $\mathrm{Ex}_C(G')$ be $\mathrm{Ex}(G')$ computed after deleting every triple that contains a far pair and has no site in $C$; $\mathrm{Ex}_C$ is monotone in $C$. Let $V(C)$ be the set of vertices of the members in $C$ and of the members incident to an edge in $C$.

\textbf{Hypothesis BC$(t)$ (bounded certificates; open).} There is $K(t,D)$ such that for every $G'$ (with the same forest $H$, and no static violation) with $\mathrm{Ex}(G')>t$, every inclusion-minimal set $C$ of sites with $\mathrm{Ex}_C(G')>t$ has $|C|\le K(t,D)$.

\textbf{Assumption (O).} $\mathrm{Ex}_C(G')$ can be computed in time $f(w)n^{O(1)}$, by a dynamic program over a tree decomposition whose bags are enlarged by the members at distance one. (Every triple with a soft pair or a cross pair has its vertices in one such bag; the argument is a sketch and is not proved here.)
\begin{theorem}[Forests, conditional]\label{thm:F}
Let $I(H)$ be a forest. Algorithm $\mathcal A$ of Appendix~\ref{app:softF} decides CEaMP; it is correct unconditionally (given the oracle). Its running time is $(6KZ^{2})^{\ell}\cdot T_{\mathrm{O}}\cdot n^{O(1)}$, where $Z=3D+3\ell+1$, $T_{\mathrm{O}}$ is the time of the oracle (O) and $K$ is the largest minimal set $C$ met. Hence, if BC$(t)$ holds for $t<\ell$ and (O) holds, CEaMP on forests is fixed-parameter tractable in $\ell$.
\end{theorem}
\emph{Proved implication}; BC and (O) are not proved. The hypothesis BC affects the \emph{running time only}, not correctness. BC$(0)$ holds with $K=2$ (Lemma~\ref{lem:sfH2}). The case $\ell\le1$ is already polynomial by Theorem~\ref{thm:A} and is not progress. Evidence for BC: minimal certificates had at most $t+2$ sites in every adversarial search ($1,3,4$ for $t=0,1,2$). A proof would need a weighted Helly statement for the five-state problem; breaks interact through locks (Proposition~\ref{prop:hub}), so an exchange argument in the style of Proposition~\ref{prop:RM2} fails.

\subsection{Certificates for far edits: a branching scheme and its unbounded branching}\label{sec:softcert}
We now use a different split of the hard pairs, tuned to a branching algorithm. For a vertex $v$ let $\mathcal M(v)$ be the set of members through $v$. A pair $\{u,v\}$ is \emph{local} if $u,v\in V(H)$ and some member through $u$ equals or shares a vertex with some member through $v$, and \emph{remote} otherwise; soft pairs are local, every pair with an end outside $V(H)$ is remote. Let $\mathcal N$ be the set of remote hard pairs. For $T\subseteq\mathcal N$ let $F'(T)$ be the least number of edits in an assignment of a status to every soft pair (each member gets a non-empty subset of its pairs) and to every local hard pair, such that every vertex triple \emph{all of whose remote pairs lie in $T$} is not an induced $P_3$ (remote pairs of $T$ keep their status in $G$; triples with a remote pair outside $T$ are dropped).
\begin{lemma}[Certificates]\label{lem:cert}
(a) $F'$ is monotone: $T\subseteq T'$ implies $F'(T)\le F'(T')$. (b) If $S$ is a solution with $S\cap T=\emptyset$, then $F'(T)\le|S|$. (c) $F'(\mathcal N)$ is the least size of a solution with no remote edit. Hence if $F'(T)>|H|+\ell$ every solution of size at most $|H|+\ell$ contains a pair of $T$ (such a $T$ is a \emph{certificate}).
\end{lemma}
\begin{theorem}[Certificate branching]\label{thm:B1}
Algorithm: if $F'(\mathcal N)\le|H|+\ell$, accept; if $\ell=0$, reject; otherwise delete pairs from $\mathcal N$ one by one while $F'>|H|+\ell$ stays true, and for every pair $e$ of the resulting minimal certificate $T$ recurse on $(G\sd e,\ell-1)$. The algorithm is correct for every packing; depth at most $\ell$ and branching factor $d=\max|T|$ over the certificates met, so $d^{\ell}(|\mathcal N|+2)$ oracle calls.
\end{theorem}
\begin{proof}
(a) fewer constraints; (b) $S$ restricted to soft and local pairs is an assignment, and every triple with its remote pairs in $T$ has them unedited; (c) with $T=\mathcal N$ nothing is dropped. Soundness: a solution of $G\sd e$ of size $\le|H|+\ell-1$ gives $S\sd\{e\}$, and $H$ is still a packing since $e$ is hard. Completeness: if no solution of size $\le|H|+\ell$ avoids all remote pairs then $F'(\mathcal N)>|H|+\ell$, a minimal certificate $T$ exists, and every solution contains some $e\in T$.
\end{proof}
So the question is how large minimal certificates are, and how $F'$ is computed. For the second, the former difficulty that a triple with two local pairs and one remote pair need not lie in a bag of a decomposition of $I(H)$ disappears for this definition of \emph{local}.
\begin{theorem}[The leaf test is fixed-parameter tractable]\label{thm:L}
Given a tree decomposition of $I(H)$ of width $w$, $F'(T)$ for any $T\subseteq\mathcal N$, in particular $F'(\mathcal N)$, is computable in time $f(w)n^{O(1)}$.
\end{theorem}
\emph{Proved} (Appendix~\ref{app:softL}; radius $3$ in $I(H)$ is attained and suffices). Remark: this is true for the \emph{adjacency-defined} notion of local used here; for locality defined through the bags of a decomposition of $I(H)$ the corresponding statement remains unproved, so the earlier gap is bypassed by changing $\mathcal N$, not closed. Together, Theorems~\ref{thm:B1} and~\ref{thm:L} give an algorithm in time $d^{\ell}f(w)n^{O(1)}$; the question is the size of $d$.

\begin{theorem}[Core bound at $\ell=0$ without soft triangles]\label{thm:CB0}
Let $H$ have no soft triangle and $\ell=0$. Every minimal $T\subseteq\mathcal N$ with $F'(T)>|H|$ has $|T|\le12$ ($\le3$ if it is a static violation), for every $c$ and $w$.
\end{theorem}
\emph{Proved} (Appendix~\ref{app:softCB}; the proof uses the far pairs as \emph{guards}, unit or two-element lists on one member, together with Theorem~\ref{thm:sfA} and a Helly lemma for $2$-element lists on a $3$-element domain). The bound $12$ is not claimed tight. At $\ell=0$ the algorithm of Theorem~\ref{thm:B1} never branches, so Theorem~\ref{thm:CB0} is a statement about the structure of $\ell=0$ certificates; its algorithmic use needs an $\ell\ge1$ version, which is open (Erd\H{o}s--P\'osa-type argument sketched, not proved). \emph{Evidence}: in $786$ random soft-triangle-free instances the largest minimal inconsistent family of guards had two members, and minimal certificates found by deletion had at most $4$ pairs.

With soft triangles the corresponding statement is false, already at $c=2$ and $w=2$.
\begin{theorem}[Unbounded minimal certificates]\label{thm:CE}
For every $N\ge1$ there is a graph $G_N$ with a pair-maximal packing $H_N$ of $3N$ members ($c=2$, $\tw(I(H_N))=2$, $N$ soft triangles, $6N+3$ vertices) and a set $T_N$ of $2N+4$ remote pairs with $F'(T_N)\ge|H_N|+1$ and $F'(T_N\setminus\{e\})=|H_N|$ for every $e\in T_N$. So $T_N$ is a minimal certificate for $\ell=0$ of unbounded size, and the minimal-certificate branching of Theorem~\ref{thm:B1} has unbounded branching factor.
\end{theorem}
\emph{Proved} (Appendix~\ref{app:softCE}; the block structure of the set of ZE vectors of $G_N$ is by computer, for the first blocks and by isomorphism of interior blocks). Two comments. (1) The proof gives $F'(T_N)\ge|H_N|+1$; the equality $F'(T_N)=|H_N|+1$ was verified by computer ($N\le10$) and is not needed. (2) At $\ell=0$ the set $T_N$ is never branched on. For $\ell=1$: the union of two disjoint copies, $T=T_N\cup T_N'$, is a minimal certificate of size $2(2N+4)$ because $F'$ is additive over disjoint unions (every triple meeting two copies has a remote pair outside $T$ and is dropped); and for \emph{connected} instances, a chain $G_{2N}$ with one further pendant $w_*$ adjacent to the glue vertex $b_N$ and the remote pairs $w_*b_N,w_*r_N,w_*p_{N+1}$ added to $T_{2N}$ gives $|T|=4N+7$ with $F'(T)-|H|=2>1$ and $F'(T\setminus\{e\})-|H|=1$ for every $e$ (evidence: $N=1,\dots,4$, exact solvers). So the branching factor of Theorem~\ref{thm:B1} with an arbitrary minimal certificate is unbounded already at $\ell=1$.

The mechanism is a chain of soft triangles that carries one bit through the glue vertices; a $2$-element guard $\{1,2\}$ (``the pair $ac$ is apart'', obtained from a pendant vertex adjacent to one end of the pair) switches off the escape state $3$ of its block. This is exactly the \emph{multi-valued unary guard} that cannot occur without a soft triangle (Theorem~\ref{thm:CB0}).
\begin{remark}[Not a hardness result; minimum certificates]\label{rem:mincert}
In $G_N$ the real optimum is $|H|+N$: every pendant edge forces its anchor to be a singleton, impossible for a centre, so all $N$ stubs are cut; the persistent bit exists only in the relaxation. Moreover there is always a small certificate there: if a far hard edge $wx$ ends at the centre $x$ of a member $(y,x,y')$ and $wy,wy'$ are far hard non-edges, then $\{wx,wy,wy'\}$ is a certificate at $\ell=0$ (both pairs at $x$ would have to be apart, two edits in one member). Every $G_N$ has such a pendant at each centre of $A_j$. So Theorem~\ref{thm:CE} refutes ``every minimal certificate is bounded'', not ``some certificate is bounded''; this is a property of the instance design, and computing a minimum certificate is itself not known to be fixed-parameter tractable. Whether a certificate of bounded size can always be chosen is open.
\end{remark}

\begin{lemma}[No conservative near-unanimity polymorphism with soft triangles, $c\ge3$]\label{lem:R5}
Let $R_k=\{111,113,122,131,212,221,222,232,311,331\}$ (states written as three-digit words in $\{1,2,3\}$). Any constraint language that contains $R_k$ and all unary relations has no conservative near-unanimity polymorphism of any arity $r\ge3$.
\end{lemma}
\emph{Proved} (Appendix~\ref{app:softR5}). The relation $R_k$ is the conjunction of all triple constraints whose scope lies in $\{A,B,C\}$ for three pairwise adjacent members $A,B,C$ on six vertices, in a packing with $c=3$ that has four further members, each owning one of the four pairs of these six vertices that are neither in $A,B,C$ nor hard non-edges (an explicit instance, found and checked by computer). It requires $c\ge3$; for $c=2$, where Theorem~\ref{thm:CE} lives, the lemma says nothing. What explains Theorem~\ref{thm:CE} at $c=2$ is that a near-unanimity operation of arity $r$ bounds the cores of unit lists but not of $2$-element lists, since it only controls coordinates where the other $r-1$ values coincide. The lemma shows that the route through a uniform near-unanimity polymorphism is closed for the full language with soft triangles; it is not a hardness proof.

\subsection{Abstract hardness and the obstruction to realising it}\label{sec:softhard}
Recall that W[1]-hardness through hard edges between far vertices needs an operation of cost one that repairs an obstruction at one of many interior positions. We separate the abstract problem from its realisation.
\begin{definition}[LPPH]
\emph{Linked-pair path hitting}: a forest, a family of \emph{obstruction paths} in it, a set of \emph{links} (pairs of nodes) and a number $L$. An operation is a single node or a link; it hits every path containing one of its nodes. Can $L$ operations hit all paths?
\end{definition}
\begin{theorem}[Abstract hardness]\label{thm:B}
LPPH is W[1]-hard parameterized by $L$, already for a forest of $2k$ paths and a set of links that is a perfect matching (every node in at most one link), with $L=k^2$.
\end{theorem}
\emph{Proved} (Appendix~\ref{app:softB}), by a reduction from Grid Tiling with $\le$~\cite{Cygan15}; cross-checked by brute force on $1{,}710$ instances (two independent solvers, no mismatch). It is the abstract form of the mechanism ``a far two-sided edit shares cost between two places''; an operation can hit an interior node of an obstruction, which is the feature the reduction uses.

\begin{theorem}[Tunnel releases cannot realise Theorem~\ref{thm:B} without soft triangles]\label{thm:C}
Let $(G,H)$ be soft-triangle-free and consider the tunnel-release model (Appendix~\ref{app:tunnels}).
\begin{enumerate}
\item[(i)] Let the \emph{conflict graph} $\mathcal C$ have as vertices the anchors that carry a tunnel or a stub, an edge $xy$ if $\{x,y\}$ is not a face, and a loop at $x$ if $\{x\}$ is not a face. Releasing a set $R$ of these anchors (cutting their tunnels) gives a feasible instance iff $R$ is a vertex cover of $\mathcal C$ (loops included); a cut at $x$ repairs exactly the conflicts at $x$, and every conflict has at most two ends.
\item[(ii)] Let $X$ be a set of anchors and $R_0$ a set of released anchors disjoint from $X$ such that the instance with release set $R_0$ is infeasible but $R_0\cup\{x\}$ is feasible for every $x\in X$. Then $|X|\le2$. In particular an obstruction that must be repaired by cutting a tunnel at \emph{any} one of $N\ge3$ anchors (as Theorem~\ref{thm:B} needs, with $N=|S_{ij}|$ positions) cannot occur.
\end{enumerate}
\end{theorem}
\begin{proof}
(i) A set of locked anchors is feasible iff it is a face, iff (Corollary~\ref{cor:sfcore}, using that a non-face contains a non-face of size at most two) it contains no loop and no edge of $\mathcal C$, i.e.\ iff the released set covers $\mathcal C$. (ii) $R_0$ is not a vertex cover, so some edge or loop $e$ of $\mathcal C$ is uncovered; for each $x\in X$ the set $R_0\cup\{x\}$ covers $e$, so $x$ is an end of $e$, and $e$ has at most two ends.
\end{proof}
\emph{Remarks (heuristic; not part of Theorem~\ref{thm:C}).} (iv) \emph{Breaks.} A double-edited member is an unpaired single-node operation, so the tile coupling between two arms is lost; the compatibility of a break with its neighbours is controlled by the same cross bits, so a break is a phase flip whose position can be restricted by the design but not a decoupler. We found no way to couple two break positions on different tracks at total cost $2$. In the five-state model no list of an outside vertex excludes the state $4$ (the break) without also excluding state $1$ or state $2$ (a finite check on the second row of the table of Observation~\ref{obs:lists}, and on all intersections of such lists), so no hard vertex makes the break of a chain member dearer than both of its single-edit states. (v) \emph{Other operations.} Inserted far pairs and cuts inside hard components of size at least three also act only through unary changes at their two ends (Lemma~\ref{lem:sfU}), and isolating a vertex of a hard $K_s$ costs $s-1$ cuts. We have no proof that no combination of operations simulates interior hits.

\subsection{What is known: the status}\label{sec:softstatus}
\begin{enumerate}
\item \emph{The question stays open}: we prove neither W[1]-hardness nor fixed-parameter tractability of CEaMP in $\ell+w$ (Open Problem~\ref{op:main}).
\item \emph{Soft-triangle-free instances reduce to one hypothesis.} Without soft triangles, zero excess is polynomial (Corollary~\ref{cor:sfpoly}), lock cores have at most two anchors (Corollary~\ref{cor:sfcore}), the far-lift model and the tunnel-release model are fixed-parameter tractable (Theorem~\ref{thm:U}, Corollary~\ref{cor:sfTR}), the leaf test of the branching scheme is fixed-parameter tractable (Theorem~\ref{thm:L}) and minimal certificates at $\ell=0$ have bounded size (Theorem~\ref{thm:CB0}). For forests the full problem is fixed-parameter tractable in $\ell$ conditionally on BC (Theorem~\ref{thm:F}); BC and the $\ell\ge1$ version of the core bound are open.
\item \emph{Hardness would need soft triangles, at least along the routes we can analyse.} The interior-hitting mechanism that Theorem~\ref{thm:B} needs is impossible with tunnel releases when there is no soft triangle (Theorem~\ref{thm:C}), the far-acting objects in the soft-triangle-free case are unit lists with cores of at most two (Lemma~\ref{lem:sfH2}), and the only mechanism that we know to produce unbounded certificates uses soft triangles (Theorem~\ref{thm:CE}). Other operations are excluded only heuristically; and Theorem~\ref{thm:CE} is not a hardness result.
\item \emph{With soft triangles at $c\le2$} (triangles of $I(H)$, $w\ge2$): nothing is proved beyond Theorem~\ref{thm:CE}; Remark~\ref{rem:tri2} gives evidence that the lock structure is still $2$-decomposable.
\end{enumerate}
